\clearpage
\VerdierPhysicalPageAnchor{0035}%
\VerdierSemanticAnchor{ch1-1-1-1-fiber-category-cont}%
したがって $\mathcal C$ と同じ対象集合をもち、その射が $\mathcal C$ の次数
$e$ の射であるような $\mathcal C$ の部分圏である。

\VerdierSemanticAnchor{ch1-1-1-1-graded-functor}
$(\mathcal C,G)$ と $(\mathcal C',G)$ を二つの $G$-次数付き圏とする。
$(\mathcal C,G)$ から $(\mathcal C',G)$ への\emph{次数付き関手}とは、
$\mathcal G$ 上のカルテジアン関手 $F$ のことである。ファイバー圏への制限で
得られる関手を $F_e\colon\mathcal C_e\to\mathcal C'_e$ と書く。
\VerdierSemanticAnchor{ch1-1-1-1-compatible-morphism}
最後に、次数付き関手の射 $m\colon F\to F'$ が関手の
$\mathcal G$-射であるとき、それは\emph{次数付けと両立する}という。

\VerdierSemanticAnchor{ch1-1-1-2}
\noindent\textbf{1.1.2.}\quad
$\mathcal C$ を $G$-次数付き圏、$g$ を $G$ の元とする。$\mathcal C$ の各対象対
$X,Y$ に対し、
\VerdierSemanticAnchor{ch1-1-1-2-homogeneous-hom}%
$X$ から $Y$ への次数 $g$ の射の集合を $\operatorname{Hom}^{g}(X,Y)$ と書く。
このとき：
\VerdierSemanticAnchor{ch1-1-1-2-eq1}
\begingroup
\makeatletter
\tagsleft@true
\let\veqno\@@leqno
\makeatother
\begin{equation}
  \operatorname{Hom}(X,Y)
  =\coprod_{g\in G}\operatorname{Hom}^{g}(X,Y).
  \tag{1.1.2.1}
\end{equation}
\endgroup

\VerdierSemanticAnchor{ch1-1-1-2-graded-hom-set}
したがって集合 $\operatorname{Hom}(X,Y)$ は $G$-次数付き集合の構造をもつ。
次数の関手性により、任意の射 $m\colon X\to Y$ と任意の射
$n\colon Y\to Z$ に対して：
\VerdierSemanticAnchor{ch1-1-1-2-eq2}
\begin{equation}
  \deg(n)\deg(m)=\deg(nm).
  \tag{1.1.2.2}
\end{equation}

\VerdierSemanticAnchor{ch1-1-1-2-fibered-criterion}
以上の式は $\mathcal C$ が $\mathcal G$ 上の圏であることを表すにすぎない。
次に、圏 $\mathcal C$ が $\mathcal G$ 上にファイバー化されていることを表そう。
そのためには、$\mathcal C$ の各対象が、指定された次数をもつ超カルテジアン射の
終域であると言えばよい [19]。ところが $\mathcal C$ の超カルテジアン射は
$\mathcal C$ の同型射にほかならないことが直ちに分かる。したがって
$\mathcal C$ がファイバー化されていることは、次の性質と同値である：

\VerdierSemanticAnchor{ch1-1-1-2-eq3}
\noindent\textnormal{(1.1.2.3)}\quad
\begin{minipage}[t]{0.78\textwidth}
\emph{$G$ の各元 $g$ と $\mathcal C$ の各対象 $X$ に対し、終域が $X$ で次数が
$g$ である同型射が存在する。}
\end{minipage}
\par\medskip

\VerdierSemanticAnchor{ch1-1-1-2-degree-preserving-functors}
いま $(\mathcal C,G)$ と $(\mathcal C',G)$ を二つの $G$-次数付き圏とし、
$F\colon\mathcal C\to\mathcal C'$ を次数付き関手とする。まず $F$ は
$\mathcal G$-関手であるから、射の次数を保つ。また $F$ は $\mathcal C$ の
カルテジアン射を $\mathcal C'$ のカルテジアン射へ送るが、$\mathcal C$ の
カルテジアン射は同型射なので、これは自動的である。ゆえに $\mathcal C$ から
$\mathcal C'$ への次数付き関手は、ちょうど\emph{次数を保つ関手}である。

\VerdierSemanticAnchor{ch1-1-1-2-degree-e-morphism}
最後に、次数付けと両立する次数付き関手の射 $m\colon F\to F'$ とは、単に次数
$e$ の関手の射、すなわち $\mathcal C$ の各対象 $X$ に対する射
\VerdierSemanticAnchor{ch1-1-1-2-component}
\[
  m(X)\colon F(X)\longrightarrow F'(X)
\]
